\subsection{A second family and a disjoint item bank}
\label{sec:family}
Two rules fixed before scoring decide what the runs show. The first reads margin inflation across the 45 PolyPythias runs from a delete-one-seed jackknife with eight degrees of freedom and passes when the lower endpoint exceeds one, since five size clusters are too few for the wild bootstrap. The second puts every bank-one item in half $A$ and every bank-two item in half $B$ on the same runs, which removes any component the halves share only through a common item set. It returns not supported when that cross-bank interval sits above one and the jackknife interval for the within-minus-cross difference in log squared inflation includes zero. It supports the explanation when that difference interval lies above zero and the cross-bank interval covers one, and we report any other pattern as undecided. The banks also differ in item count and exemplar overlap, and bank two draws most of its items from train splits, which pretraining data are more likely to contain, so any covariance these differences carry also enters the cross-bank estimate.

We simulated both rules after fixing them, drawing seed effects and item noise from the DataDecide margin fit (Appendix~\ref{app:transport}). Over 4,000 replicates the first rule passes in 0.259 at a true 1.244 and in 0.016 at independence, with the middle 90\% of estimates between 1.026 and 1.456. At a true 1.244 with no shared item component the second returns not supported in 0.226, supported in 0.013 and undecided in 0.761, and under independence with a 0.124 item share it returns supported in 0.033 and undecided in 0.957, so an undecided verdict is expected under either explanation and not supported doesn't exclude a shared component. A first-rule pass confirms the effect in this family. In an exploratory reading we added after fixing that rule, a failing estimate inside the simulated range is a bound and one below 1.026 counts against transfer.

\outcome{\Ronecase}{1}{R1 pass}{The second family replicates the margin result. Margin inflation across the 45 runs is \pending{R1 lambda} with jackknife interval \ci{\pending{lo}}{\pending{hi}}, against 1.244 on DataDecide, and accuracy gives \pending{acc lambda} with interval \ci{\pending{lo}}{\pending{hi}}. Each PolyPythias seed reruns one recipe, which leaves no counterpart to the batch component behind our auxiliary contrast. The per-size jackknives run from \pending{min} at \pending{size} to \pending{max} at \pending{size} (Appendix~\ref{app:transport}). Without BoolQ margin inflation is \pending{noBoolQ lambda} with interval \ci{\pending{lo}}{\pending{hi}}, \pending{which moves in the same direction as on DataDecide / which doesn't move the way it does on DataDecide}.}

\outcome{\Ronecase}{2}{R1 fail}{The second family doesn't replicate the margin result at the rule's threshold. Margin inflation across the 45 runs is \pending{R1 lambda} with jackknife interval \ci{\pending{lo}}{\pending{hi}}, which includes one, and accuracy gives \pending{acc lambda} with interval \ci{\pending{lo}}{\pending{hi}}. Five sizes with nine seeds each give 40 within-configuration contrasts against 250 in DataDecide, and under the reading above this estimate \pending{bounds the effect in this family / counts against transfer}. The per-size jackknives run from \pending{min} at \pending{size} to \pending{max} at \pending{size} (Appendix~\ref{app:transport}).}

\outcome{\Rtwocase}{1}{R2 not supported}{The shared-item explanation fails its rule in this family. With bank one as half $A$ and bank two as half $B$ on the same 45 runs, margin inflation is \pending{cross lambda} with interval \ci{\pending{lo}}{\pending{hi}}, and the within-minus-cross difference in log squared inflation has interval \ci{\pending{lo}}{\pending{hi}}, which includes zero. The check gives no evidence that a component the halves of one item set share accounts for the inflation on PolyPythias, and the upper limit of that difference, \pending{hi}, bounds such a component.}

\outcome{\Rtwocase}{2}{R2 supported}{The shared-item explanation passes its rule in this family. With bank one as half $A$ and bank two as half $B$ on the same 45 runs, margin inflation is \pending{cross lambda} with interval \ci{\pending{lo}}{\pending{hi}}, which covers one, and the within-minus-cross difference in log squared inflation has interval \ci{\pending{lo}}{\pending{hi}}, which lies above zero. At least part of the within-bank inflation in PolyPythias comes from a component the two halves of one item set share, and the cross-bank value of \pending{cross lambda} is the one we'd report for its run covariance. We can't say whether DataDecide carries the same component, because its cross-bank run needs GPU time we didn't have.}

\outcome{\Rtwocase}{3}{R2 undecided}{The identification check is undecided under its rule. With bank one as half $A$ and bank two as half $B$ on the same 45 runs, margin inflation is \pending{cross lambda} with interval \ci{\pending{lo}}{\pending{hi}}, and the within-minus-cross difference in log squared inflation has interval \ci{\pending{lo}}{\pending{hi}}. Bank two holds under a fifth of bank one's items, which widens both intervals.}

\subsection{Concentration by scoring format}
Margin correlations track scoring format (Section~\ref{sec:composition}), and cross-format off-diagonal covariance alone gives diagnostic inflation of 0.921 with interval \ci{0.801}{1.035} against 1.244 for the full matrix (Appendix~\ref{app:supplementary}). A masked matrix needn't stay positive semidefinite, and format and task content aren't separable here. Setting within-format correlations to their mean of 0.624 and cross-format ones to 0.007, with the benchmark seed variances kept, gives 1.285 against the observed 1.244 and 1.728 against 1.786 without BoolQ. With equal variances the same correlations would give 1.803, a gap mostly due to BoolQ's 68\% trace share. The fit is in sample, and Appendix~\ref{app:formatfit} checks it with each benchmark left out.

\subsection{Predicting aggregate uncertainty}
\label{sec:prediction}
We find covariance gains on margins and no separation on accuracy. We compared five covariance models and an exploratory rescaled plug-in on five recipe folds, fitting the seed-covariance models off fold and giving them the held-out marginal seed standard deviations (Table~\ref{tab:prediction} in Appendix~\ref{app:supplementary}). On margins the rank-one, rank-one plus gain, and full models lower mean squared error of log aggregate standard deviation by about 0.03 from the 0.0358 of independence, and their paired recipe-bootstrap intervals exclude zero. However, the rank-one interval reaches zero when nonpositive held-out variances are filled from training folds. On accuracy every such interval includes zero, and an operational comparison refitting without held-out marginals separates no model from independence.

Over the 1,445 same-size recipe pairs whose margin paired difference has a positive variance estimate, its standard deviation has median ratio 1.024 to the independence value (Appendix~\ref{app:paired}). On real comparisons at each size below 1B, we called a pair of recipes when their aggregate gap exceeded two standard errors, computed either under independence or with the estimated covariance, leaving the pair out of both estimates. We scored each call against the pair's ordering at 1B whenever that ordering itself cleared two corrected standard errors. Across the four sizes and both scales, the correction removes between 1 and 12 of the 173 to 223 calls on such pairs, which changes the share of wrong calls by at most 0.009, from rates between 0.029 and 0.070, and every recipe-bootstrap interval for that change reaches zero (Table~\ref{tab:decision} in Appendix~\ref{app:decision}). A wrong call here also counts genuine rank changes between sizes, so these rates aren't false-positive rates in the usual sense. In single-run comparisons the correction changes the wrong-call rate by at most 0.019 for runs sharing a seed label and 0.016 for runs that don't, every interval reaching zero, and the two kinds of pair have paired-difference deviations within a factor of 0.928 to 1.025, so seed matching doesn't explain the few changed calls.


\section{Related work}
\citet{bouthillier2021} derive variance inflation from correlated trainings on one task, whereas we account for covariance across benchmarks within a replicate run. \citet{reimers2017}, \citet{zhou2020}, \citet{madaan2024}, and \citet{fehlauer2025} study seed-related variability, and \citet{dodge2020} test whether good initialisations transfer across tasks. \citet{jordan2024} finds run deviations in accuracy nearly independent across ImageNet evaluation sets, which our accuracy interval allows, while our margin covariance sits within formats. \citet{miller2024} gives clustered and paired item-sampling errors with run variance fixed, and \citet{henderson2018} and \citet{agarwal2021} resample runs within a task, which discards the paired run effects we keep. We know of no study that estimates the covariance of run noise between benchmarks or an effective benchmark count from that covariance.

Our cross-half estimator adapts the split-half construction of \citet{spearman1910,brown1910}, the attenuation correction of \citet{spearman1904}, and the variance components of \citet{searle1992} and generalizability theory \citep{brennan2001}, and what we add is the estimand and its measurement. \citet{cheverud1988} motivates our phenotypic plug-in comparison (Table~\ref{tab:prediction}), and we borrow the covariance analogy of \citet{martin1977} without its variance components.

Our few-cluster inference follows \citet{cameron2008} and \citet{mackinnon2017}. \citet{heineman2025} relate per-benchmark signal and noise on DataDecide and OLMo checkpoints to decision accuracy without estimating covariance between benchmarks. Across models, benchmark scores share a low-dimensional capability structure \citep{ruan2024,burnell2023revealing}, whereas we measure covariance of run deviations with the configuration held fixed.

\section{Discussion and limitations}
We read the accuracy result as a bound rather than as evidence of independence. The full-battery accuracy interval includes one at accuracy-like item noise, a design that would detect a true ratio of 1.10 in 0.135 of replicates, and the primary analysis gives an upper limit of 1.157 at the scored checkpoint. Exploratory cuts whose intervals exclude one, without BoolQ at 1.558, at the adjacent checkpoint and on the held-out tasks, don't settle it.

Apart from the registered test, which failed at its registered scope, and the two PolyPythias rules fixed before scoring, our study remains exploratory. We didn't score the 375 DataDecide runs on the second bank, which leaves the 1.244 untested on fresh items and the shared-item check unrun in the family that supplies our headline. Our populations stop at 1B and at likelihood-scored multiple choice, since no larger public multi-seed population shares replicate item scores (Appendix~\ref{app:design}) and generative tasks lack our margin. A scalar item effect shared across benchmarks and repeated across halves would enter the numerator whole, and shares of 0.124 and 0.009 of item-noise variance reproduce 1.244 and 1.078, while at a share of 0.05 only 0.007 of simulated replicates reach the margin value. We haven't measured such an effect or separated it from seed covariance (Appendix~\ref{app:calibration}).\outcome{\Rtwocase}{1}{R2 not supported}{ The cross-bank check in PolyPythias constrains that item effect in one family, and it can't reach an item effect that the two banks share.}\outcome{\Rtwocase}{2}{R2 supported}{ The cross-bank check in PolyPythias finds such a component in one family, so the within-bank value can overstate run covariance there.}\outcome{\Rtwocase}{3}{R2 undecided}{ The cross-bank check in PolyPythias leaves that item effect undecided in one family.} A scalar component shared by every benchmark would also enter the cross-format pairs, whose diagnostic of 0.921 constrains a uniform component but not one that follows scoring format.

We cluster on recipes, but the release shares one seed label across every recipe inside a size band, and a run effect following the band would break that choice and take simulated margin coverage over 10,000 replicates to 0.823 at a quarter of latent variance and 0.581 at a half, although the size-band-clustered margin interval \ci{1.030}{1.426} of Table~\ref{tab:primary} still excludes one. A permutation test that aligns runs by batch slot bounds that effect at 0.044 of the seed covariance on margins, where simulated coverage is 0.941 at a share of 0.05, and that test's power is 0.490 at the bound. At 1B alone the cluster-$t$ margin interval of 0.811 to 1.628 includes one (Appendix~\ref{app:supplementary}).

We recommend paired replicate scores on the chosen battery, which give the aggregate seed standard deviation directly, and a benchmark-removal check. Against one configuration's six-run margin ratio on the earlier three-size PolyPythias runs, whose middle 90\% runs from 0.70 to 1.85, independence predicts with mean absolute log error 0.339 against 0.536 for the plug-in and 0.512 for our decomposition, while against a target pooled across configurations the decomposition beats independence, 0.275 against 0.319 on margins and 0.367 against 0.523 on accuracy (Appendix~\ref{app:transport}). Our ratio therefore measures this battery and helps only where replicates are missing and the target battery's own ratio is close. Without replicates, the format fit of Appendix~\ref{app:formatfit} gives a fallback scaling we haven't tested outside this battery. Scaling the independence standard error by a ratio from the same margin population returns the simulated false-positive rate from 0.113 to 0.051, and a separate transfer simulation to a battery whose own ratio is 1.5 leaves it at 0.101. In a population with the released item counts and equicorrelated seed covariance, three runs each let the margin interval exclude one in 0.997 of replicates, and accuracy needs eight to reach 0.803.

\label{maintext:end}
\section*{AI use statement}
We used generative AI tools to propose and refine hypotheses, critique the estimator and simulation design, assist with analysis code, revise manuscript text, survey literature, and contribute to the technical and editorial review of this version. The study includes synthetic calibration and sensitivity simulations, whose role differs from the released model outputs used for the empirical analysis. The authors remain responsible for verifying AI-assisted claims, code, citations, and the final submission, including an accurate disclosure of AI involvement in generating simulation data.

\section*{Ethics statement}
This study analyses released model outputs and introduces no new participant recruitment or data collection. Its intended use is to improve uncertainty reporting in model evaluation. An independence approximation used outside the studied battery could understate uncertainty, which keeps the recommendations conditional on the measured population and score definition. Use of derived artifacts must retain the applicable source attribution and licensing requirements.

\section*{Reproducibility statement}
Section~\ref{sec:estimator} and Appendix~\ref{app:derivations} define the estimator and its assumptions. Appendix~\ref{app:implementation} specifies the interval calculation and data reduction. Appendix~\ref{app:design} records the populations and planned checks. Appendix~\ref{app:supplementary} reports later sensitivity analyses alongside the primary results. The released archive also holds an extended supplement with the full-length versions of the compressed appendix sections. The primary calculation uses 4,999 bootstrap draws, 2,000 calibration replicates, 500 permutation replicates, 50 item re-splits, and master seed 20260101, and reproduction requires the per-item reductions, item ordering and half assignments, selected checkpoint identifiers, and the corresponding analysis scripts. We release all four of them alongside the paper. Neither the analysis plan nor the protocol and rule commits carry an independently corroborated timestamp, and the plan's screening split was never formed. A separate implementation run on Kaggle reproduces the primary estimates, the BoolQ-removal estimates, and their intervals to three decimals from the released per-item scores, a run that also produced the passage-aware split, the coverage simulation, the operational prediction comparison, and the seed-matched pair analysis. The PolyPythias transfer, five-size rescore, and checkpoint results come from earlier execution records that we didn't rerun for the final analysis. The registered test ships with its protocol, the dated amendments that record each change to it, the frozen hash of its request file, and the kernels that scored and analysed it, and the last twenty runs at 1B were scored off Kaggle on three RTX 3090 cards. The PolyPythias battery and the second bank were scored in float32 with transformers 4.57.1 on Kaggle T4 cards, each run checking its cached scores against an uncached reference on 24 items to within 0.001 nats, and the release ships the frozen hashes of both request files with the kernels that scored them. Separate CPU kernels on the released runs produced the recipe-decision, seed-count, battery-size, and item-count analyses.
